We use the following basic task management primitives to
develop the Lean frontend.
\begin{lstlisting}
Task.mk :$\ $ (Unit -> $\alpha$) -> Task $\alpha$
Task.bind : Task $\alpha$ -> ($\alpha$ -> Task $\beta$) -> Task $\beta$
Task.get : Task $\alpha$ -> $\alpha$
\end{lstlisting}
The function $\textit{Task.mk}$ converts a closure into a \kw{task} value
and executes it in a separate thread,
$\textit{Task.bind}\ t\ f$ creates a \kw{task} value that waits for $t$ to finish and
produce result $a$, and then starts $f\ a$ and waits for it to finish.
Finally, $\textit{Task.get}\ t$ waits for $t$ to finish and returns the
value produced by it. These primitives are part of the Lean runtime,
implemented in C++, and are available to regular users.

The standard way of implementing thread safe reference counting uses
memory fences~\cite{Boost}.  The reference counters are incremented
using an atomic fetch and add operation with a relaxed memory
order. The relaxed memory order can be used because new references to
a value can only be formed from an existing reference, and passing
an existing reference from one thread to another must already provide
any required synchronization. When decrementing a reference counter, it
is important to enforce that any decrements of the counter from other threads
are visible before checking if the object should be deleted. The
standard way of achieving this effect uses a \emph{release} operation
after dropping a reference, and an \emph{acquire} operation before
the deletion check. This approach has been used in the previous
version of the Lean compiler, and we have observed that the memory
fences have a significant performance impact even when only one thread
is being executed. This is quite unfortunate because most values
are only touched by a single execution thread.

We have addressed this performance problem in our runtime by tagging
values as \emph{single-threaded}, \emph{multi-threaded}, or
\emph{persistent}. As the name suggests, a single-threaded value is
accessed by a single thread and a multi-threaded one by one or more
threads. If a value is tagged as single-threaded, we do not use any
memory fence for incrementing or decrementing its reference counter.
Persistent values are never deallocated and do not even need a
reference counter.  We use persistent values to implement values
that are created at program initialization time and remain alive
until program termination.  Our runtime enforces the following
invariant: from persistent values, we can only reach other persistent
values, and from multi-threaded values, we can only reach persistent
or multi-threaded values. There are no constraints on the kind of
value that can be reached from a single-threaded value. By default,
values are single-threaded, and our runtime provides a
$\textit{markMT}(o)$ procedure that tags all single-threaded values
reachable from $o$ as multi-threaded.  This procedure is used to
implement $\textit{Task.mk}\ f$ and $\textit{Task.bind x f}$.  We use
$\textit{markMT}(f)$ and $\textit{markMT}(x)$ to ensure that all
values reachable from these values are tagged as multi-threaded \emph{before} we
create a new task, that is, while they are still accessible from only one thread.
Our invariant ensures that $\textit{markMT}$ does
not need to visit values reachable from a value already tagged
as multi-threaded.  Thus values are visited at most once by
$\textit{markMT}$ during program execution. Note that task creation is not
a constant time operation in our approach because it is proportional to the number of
single-threaded values reachable from $x$ and $f$. This does not seem
to be a problem in practice, but if it becomes an issue we can provide a
primitive $\textit{asMT}\ g$ that ensures that all values allocated
when executing $g$ are immediately tagged as multi-threaded.
Users would then use this flag in code that creates the values reachable by
$\textit{Task.mk}\ f$ and $\textit{Task.bind x f}$.

The reference counting operations perform an extra operation to test
the value tag and decide whether a memory fence is needed or not.
This additional test does not require any synchronization because the
tag is only modified before a value is shared with other execution
threads. In the experimental section, we demonstrate that this simple
approach significantly boosts performance. This is not surprising
because the additional test is much cheaper than memory fences on
modern hardware.
The approach above can be adapted to more complex libraries for
writing multi-threaded code. We just need to identify which functions
may send values to other execution threads, and use $\textit{markMT}$.

%%% Local Variables:
%%% mode: latex
%%% TeX-master: "refcount"
%%% End:
